\documentclass[10pt]{article}
\usepackage[margin=1in]{geometry}
\usepackage[T1]{fontenc}
\usepackage{lmodern,microtype,amsmath,amssymb,amsthm}
\theoremstyle{definition}
\newtheorem{assumption}{Assumption}
\begin{document}
\section*{Symmetric Gaussian Reverse-Fisher Recovery}
We consider recovering a hidden vector $\mu\in\mathbb{R}^d$, up to sign, using gradient descent on the reverse-Fisher loss between symmetric Gaussian mixtures with identity covariance. The public parameters are the dimension $d\geq 1$, norm bounds $0\leq a\leq b$ satisfying $a\leq\|\mu\|\leq b$, accuracy $\varepsilon\in(0,1]$, and failure tolerance $\delta\in(0,1]$. Each oracle call returns the exact loss gradient at a queried vector. Updates use strictly positive step sizes, and the output is the current iterate. The required measurable output $\bar{x}\in\mathbb{R}^d$ satisfies $\min\{\|\bar{x}-\mu\|^2,\|\bar{x}+\mu\|^2\}\leq\varepsilon$ with probability at least $1-\delta$.

Equip $\mathbb{R}^d$ with its Euclidean inner product and norm. For a candidate vector $x\in\mathbb{R}^d$, define the mixture and its score by
\[
\begin{aligned}
q_x &= \tfrac12\mathcal{N}(x,I_d)+\tfrac12\mathcal{N}(-x,I_d),\\
s_x(y) &= \tanh(\langle x,y\rangle)x-y,
\qquad y\in\mathbb{R}^d,
\end{aligned}
\]
where $I_d$ is the identity matrix. Writing $\gamma_d=\mathcal{N}(0,I_d)$, the optimization objective is
\[
\min_{x\in\mathbb{R}^d} F_\mu(x),
\]
with the reverse-Fisher loss defined explicitly as
\[
\begin{aligned}
F_\mu(x)
&=\frac12\int_{\mathbb{R}^d}
 \|s_x(x+z)-s_\mu(x+z)\|^2\,\gamma_d(dz)\\
&\quad+\frac12\int_{\mathbb{R}^d}
 \|s_x(-x+z)-s_\mu(-x+z)\|^2\,\gamma_d(dz).
\end{aligned}
\]
Thus the score discrepancy is integrated under the candidate mixture $q_x$. The admissible instances, oracle access, and algorithm family are specified below.

\begin{assumption}[Parameters and hidden target]
The parameters satisfy
\[
\begin{aligned}
d&\in\{1,2,\ldots\}, & 0&\leq a\leq b,\\
0&<\varepsilon\leq1, & 0&<\delta\leq1.
\end{aligned}
\]
The target is a vector $\mu\in\mathbb{R}^d$ with $a\leq\|\mu\|\leq b$. The full tuple $(d,a,b,\varepsilon,\delta)$ is available to the algorithm; $\mu$ is hidden.
\end{assumption}

\begin{assumption}[Exact-gradient oracle]
Oracle ID $0$ accepts any query $x\in\mathbb{R}^d$ and returns $\nabla F_\mu(x)\in\mathbb{R}^d$. This identity holds at every call index and for every realization of oracle randomness. Each query costs one oracle call. Internal algorithm randomness is independent of oracle randomness.
\end{assumption}

\begin{assumption}[Gradient-descent family]
The algorithm chooses an initial point $x^0\in\mathbb{R}^d$. At each iteration $k$, it queries its current point $x^k$, chooses a strictly positive step size $\eta_k$ before receiving the reply, and updates according to
\[
x^{k+1}=x^k-\eta_k\nabla F_\mu(x^k),
\qquad \eta_k>0.
\]
The initialization and step choices may depend on the public parameters and internal randomness, and each step choice may also depend on the preceding oracle history. The returned point is always the current iterate.
\end{assumption}

An externally specified budget $K\in\mathbb{N}$ may depend on $(d,a,b,\varepsilon,\delta)$ and counts calls to oracle ID $0$. After $K$ calls, the output is $\bar{x}=x^K$. Let $\mathbb{P}$ denote the joint law of algorithm and oracle randomness. Success requires that $\bar{x}$ be measurable and that
\[
\mathbb{P}\!\left(
\min\{\|\bar{x}-\mu\|^2,\|\bar{x}+\mu\|^2\}
\leq\varepsilon
\right)\geq1-\delta.
\]
\end{document}
